\section{Case Study: Polymorphism in a multi-language Ray-Tracer}\label{s:implementation}

In the previous sections, we discussed some of the theoretical underpinnings of language interoperability and their relationship to the practical aspects of implementing cross-language tooling. A key emphasis of our exploration, especially in the latter half, was comparing and contrasting the attributes of interoperability tooling between similar and dissimilar languages and their respective runtime models. The idea of reconciling the semantics of different languages and leveraging their respective strengths is, in large part, what motivated the exploration of this topic. Julia, especially, is an interesting candidate for the case study, as its performance characteristics are considerably more dependent on an idiomatic programming style than those of most other languages. Idiomatically written code, regardless of the language, tends to be more performant, as the set of paradigms it applies tends to enable a more effective use of the underlying execution model. Phrased differently, we can understand idiomatic code to be defined by a set of soft or non-strictly enforced requirements as defined by the language execution model. Julia's execution model, the JIT compiler, imposes a more nuanced and impactful set of requirements on idiomatic code, which in turn can lead to a higher learning curve for new users. 

One of Julia's main selling points is that \textit{when written well}, it has performance characteristics close to that of C while retaining the high-level ease of use of Python. The somewhat unique nature of Julia, or more specifically, its JIT compiler, is what served as the primary motivation to use a ray-tracer implementation written in Julia as the main testbench for this exploration. The ray-tracer itself is a relatively simple implementation based in large parts on Peter Shirley's guide \cite{Shirley2024RTW1, Shirley2024RTW2} and his C++ implementation. The initial Julia implementation exhibited performance characteristics that were less than ideal when compared to Peter Shirley's equivalent C++ implementation. To elaborate on why the non-idiomatic nature of the implementation led to such stark performance differences, we will briefly explain how Julia's JIT compiler works. A quote from the Julia forums summarizes this nicely. 

\begin{quote}
  \textit{... in short, it's a reasonably good static compiler but a terrible JIT}
\end{quote}

The first response some people might have to this is confusion. Julia is a JIT-compiled language, so how is its JIT compiler terrible. Well, this comes down in large part to how its JIT compilation works. The core of this is how it performs online partial evaluation (OPE) concerning its type system. OPE w.r.t, the type system in this instance, refers to the process of, at runtime (online), Julia specializes (partially evaluates) methods for which it can infer the types of every local variable from the types of the arguments or any other constants; in this instance, a method is said to be `type stable'. Specializing type stable methods means it compiles methods for the set of concrete types it can infer. Practically speaking, this mechanism can also be viewed as a form of monomorphizing a set of functions for which the JIT compiler can determine the concrete types, thereby caching their compiled instances for fast execution. The reason, then, why Julia's JIT compiler could be called `terrible' would be that it, currently at least, does not implement profile-guided optimization (PGO). The lack of PGO and its reliance on type-stable code for speed means that there is a strong correlation between idiomatically written Julia code and its performance. 

Naturally, there are various approaches one can take to deal with a situation like this. First and foremost, this would generally involve writing the code in a type-stable manner, which in turn would optimize performance by reducing the number of cache misses on precompiled methods. However, a case could also be made for simply using existing high-performance code, which avoids rewriting altogether. In practice, using existing high-performance libraries in dynamic languages, especially in the context of scientific computing, usually means using libraries that wrap, often via FFI, fast, low-level frameworks like those that implement the Basic Linear Algebra Subprograms (BLAS) specification for everyday linear algebra operations. An idea that arose from this was to replace the most performance-intensive component of the ray tracer entirely with implementations in C++ and Python, rather than simply using existing libraries to replace slower operations. The idea here is twofold: firstly, we can assess the viability and performance characteristics of interoperability when applied to a more complex system. Secondly, we can also compare and contrast the earlier discussed relations between language similarity and overhead.    

\subsection{Project Context \& Toolchain}


As one of our objectives is not only to research the theoretical underpinnings of interoperability but also to assess their viability in practice, we now describe how many of the previously discussed topics are relevant in real-world scenarios. In the preamble, we decided to use a ray-tracer implementation written in Julia as the main subject for this exploration. Before delving into the implementation of cross-language aspects, it's worthwhile to provide some motivation for why a ray tracer serves as a suitable subject. Although the decision was initially ad hoc, it presents several potentially insightful attributes. Ray-tracers, especially more developed legacy implementations, are well-known for their complexity—not only due to their foundation in mathematical modeling but also because of the interconnectedness of various subsystems required for image generation. Even minimal ray-tracer implementations present a situation analogous to more complex systems, offering broader insights into the effects of interoperability compared to domain-specific minimal examples. 

Performance is another key aspect of ray tracers. While their origins lie in drawing simple images, modern ray tracers—or, more broadly, physically based rendering (PBR) frameworks—are used in real-time in video games, creating a significant demand for high-performance PBR systems. Although the ray-tracer used in this paper implements only a minimal subset, it directly relates to a domain where performance is critical, serving as an indicator of the viability of the applied interoperability techniques in performance-demanding scenarios. Finally, despite the complexity of ray tracers, they tend to be modular systems in certain respects. This modularity allows us to more easily separate logically distinct components, making the rewriting process more manageable. A high-level overview of the ray tracer can be seen in Figure \ref{fig:simple_rt_impl}.

\begin{figure}[ht!]
  \centering
  \begin{tikzpicture}[subgraph text bottom=text centered, subgraph nodes={font=\itshape, black}, nodes=draw]
    \graph [tree layout, grow=down, sibling sep=5mm, level sep=5mm, component align=center, component direction=270] {
        "Image generation" <-> {
            "Hit detection" -> [loop below] "Hit detection",
            "Scattering"
        },
    };
  \end{tikzpicture}
  \caption{Ray-Tracer high-level control flow}
  \label{fig:simple_rt_impl}
  \Description{}
\end{figure}

\subsubsection{Implementation approach}

Now that we have established our motivation for testing and have an implementation to work with, the next step is figuring out how to test it in practice. To isolate the most performance-impacting component, we profile the application, analyzing both runtime and memory usage before making informed judgments based on this data. Regarding language choices for rewriting, our constraints are those with embedding APIs for Julia, but we still have several options. It makes sense to choose two languages with fundamentally different properties—Python and C++—as they provide a broad spectrum of insights applicable to similar domains. Both languages have well-established interface libraries, making them practical choices for evaluation. To test the effectiveness of the component, we use benchmarking tools to analyze its performance in isolation, along with Julia-sided tools to assess the combined performance of the component, interface code, and any overhead introduced.

\subsubsection{Choice of Binding Generators}

Writing the components alone is not enough if we can't integrate them into the larger system, so the next step was selecting an embedding API to attach the hit functions to the ray tracer implementation. For Python, the choice was straightforward—Python/JuliaCall \cite{Rowley2025JuliaPy} was used, which builds on an earlier implementation and supports a broader range of conversions between the two languages. C++, however, presented a more interesting challenge. Unlike Python, which has trivial reflection mechanisms making object conversion straightforward, C++ requires more deliberate choices. The two main options were CxxWrap \cite{Janssens2025JuliaInterop}, which has a strong feature set but lacks user-friendliness, and Jluna \cite{jluna}, a newer and more user-friendly alternative that is missing some mapping mechanisms. We chose Jluna because the focus was on minimizing additional code and effort while creating usable components in another language. Jluna, leveraging modern and extensible C++ techniques, proved to be more future-proof and developer-friendly. We can tie the future outlook to a broader consideration: when rewriting a component in another language, if the results are positive, ideally, we want minimal refactoring of our existing codebase to integrate it into the larger system. This means choosing an API that (1) requires the least additional code and (2) minimizes the code needed to attach the component. One key takeaway from this process, and something which might already be quite evident given the discussion in \ref{s:background}, is that working between Python and Julia already poses significantly fewer challenges when searching for a binding generator due to the added feature set provided by both languages being designed with dynamism—and by extension runtime reflection—in mind.

\subsubsection{Implementation setup}

The code for setting up the scene, which includes parsing the `.obj' files that describe the vertices that bound all faces in the scene, loading the faces into a bounding volume hierarchy (BVH), and managing the actual image generation by drawing the pixels, is all uniformly done in Julia. The isolated components, namely the tree search and hit computation for a given ray cast, are rewritten in C++ and Python and then replaced in the Julia implementation, depending on which language is used to initiate the ray tracer. In both instances, we redefine the hit functions such that the Julia code is redirected to the language-specific variants, which act as proxies to the actual host language code.

\subsection{Design Challenges}

\subsubsection{Problem Statement}

Now, while the implementation here was straightforward in theory, a significant issue encountered was related to the constraints imposed by the binding generator between Julia and C++. To explain the cause of this issue, it's helpful first to outline what the implementation needs to do, or more specifically, what we are moving across the language boundary. As the BVH traversal served as the most logical component to isolate due to its aforementioned separation from other aspects, we chose it to serve as the bridge. The three main things needed in our implementation when trying to do hit detection are the BVH, which spatially structures the objects in the scene; the actual ray for which we are computing the hit; and finally, a hit record, which stores the information from the hit (e.g., material, normal vector to the surface, etc.). Assessing the complexity of these types, the simplest one is the ray, as it just stores an origin and direction. Next is the hit record, which has polymorphic attributes, namely the material of the object that the ray hit. Finally, there is the BVH tree, which, as implied in the name, is a tree structure that adds another layer of complexity, as there are now recursive polymorphic attributes with the left and right child being BVH nodes, which are subtypes that define the underlying type of the shape represented (e.g., Triangle, Sphere, etc.). 

In both instances, the type of polymorphism being discussed is subtyping polymorphism. While both C++ and Julia, abstractly speaking, implement similar conventions for enabling subtyping polymorphism—namely, some abstract type and a means of specifying a nominal type hierarchy—there are considerable differences when it comes to the practical case of interoperability. A key distinction between the high-level notions of interoperability and the low-level practicalities is how types are represented. Since we are fundamentally working through the C-ABI, or more accurately, through an abstraction (Jluna) over the Julia C-API, before we do anything, we need to parse the underlying representation of Julia types. Fortunately, this process is elegantly implemented within Jluna for all types, including structs, by the original author. Jluna achieves this mainly by using generics to deconstruct the low-level representation of a Julia type recursively. To ensure constraints on the implementation—by which we mean rejecting types that cannot be parsed for some C++ analog—the implementation utilizes constraints (concepts, in this case) on its generic deconstruction methods, which ensure certain invariants hold that constrain the set of bridgeable types to only those specified by the programmer. One of these constraints was a static assertion that prevented bridging abstract types—in other words, types that exhibited subtyping polymorphism—which brings us back to the issue we encountered. The question then becomes: How do we solve this problem? Trivially, a straightforward solution to the C++ and Julia interoperability would have been to rewrite the ray tracer such that the BVH and hit record represent everything in a non-polymorphic format. However, there are a few reasons we were motivation to bridge these types.

\begin{enumerate}
  \item \textit{Challenging Common Notions:} it's often assumed in the context of interoperability for larger organizations that the set of bridgeable behaviors should be restricted to minimal constructs. We sought to challenge this assumption; as type systems in mainstream languages become increasingly expressive, a highly restrictive approach to cross-language systems may eventually limit the applicability of libraries that leverage more advanced type system features.
  \item \textit{Flexibility Discrepancy:} Additionally, we aimed to clarify whether the observed discrepancy in flexibility between dynamic and statically compiled languages in cross-language systems is inherent to the nature of type systems and execution models or if it's primarily a consequence of specific implementation philosophies or limited use of language features.
  \item \textit{Viability of Loosening Constraints:} We were interested in evaluating the viability of relaxing constraints, both in terms of the resulting implementation complexity and the performance implications associated with bridging such types.
  \item \textit{Feasibility Exploration:} Finally, we aimed to explore the overall feasibility of complex type bridging, given its relative rarity in cross-language systems involving statically typed languages—particularly when the statically typed language is the primary focus rather than the dynamic language. We anticipated that this exploration could yield valuable insights.
\end{enumerate}

To contrast this with the case of Python and Julia, the story was interesting in its own right. JuliaCall, the library we chose for bridging between these two languages, wraps Julia types with opaque \texttt{juliacall} class objects that idiomatically mirror their Julia counterparts in terms of field access. The use of opaque wrappers, in turn, means that unless we implement custom conversion rules to produce the actual class objects corresponding to those we defined within the Python implementation, we don't have any subtyping relationship. We can still emulate nominal subtyping by overloading pythons implementation of the singledispatch decorator using unique properties the opaque wrapper provides, though this is a rather cumbersome solution. 

With the broad idea in mind of what we want to achieve now and some motivation as to why, it makes sense to establish a more concrete approach. Before we assess the details for how to extend the Jluna library, we should first assess in detail how the library works currently. As discussed throughout \ref{s:background}, if we think of C++ as our host language and Julia as the guest, then boxing would be the operation of representing a C++ type in Julia, and unboxing would be the reverse operation. In Jluna, boxing and unboxing have multiple steps; we can give a high-level overview of it as follows, with a general breakdown of what each step entails.

\begin{figure}[ht!]
  \centering
  \begin{tikzpicture}[subgraph text top=text centered, subgraph nodes={font=\itshape, black}]
  \graph [layered layout, grow=down, components go right, sibling sep=5mm, level sep=2.5mm, component align=center] {
    "\textbf{C++}"   ,
    // [component direction=270, nodes=draw] {
      "compile time"[draw] // [layered layout, grow=down] {
        "type specification",
        "type implementation"
      },
      "type specification" -> "type implementation",
      "runtime"[draw] // [layered layout, grow=down] {
        "unboxing",
        "boxing"
      },
      % Connect compile time to runtime
      "compile time" -> "runtime"
    },
    "\textbf{Julia}"
  };
  % Add external arrows for boxing/unboxing
  \draw[->] (boxing) -- +(1,0);
  \draw[-] (boxing) -- +(-1,0);
  \draw[->] (unboxing) -- +(-1,0);
  \draw[-] (unboxing) -- +(1,0);
\end{tikzpicture}
  \caption{High-level overview of Jluna application}
  \label{fig:jluna_high_level_impl}
  \Description{}
\end{figure}

\begin{enumerate}
  \item \textit{Type specification}: For the type specification of object types (i.e., classes and structs), the programmer defines a set of getter and setter functions in addition in addition to some other type attributes, which are then encoded in a meta-type akin to meta-types generated via dynamic reflection mechanisms in dynamic languages.
  \item \textit{Type implementation}: Once this type is defined, the programmer calls an implementation method that utilizes the meta-type to generate an idiomatically mirrored Julia type using its metaprogramming features, which creates the corresponding type within the Julia runtime.
  \item \textit{Unboxing}: Unboxing is then applied to all arguments of methods that are exposed as cross-language via an implicit unboxing proxy, which is generated through the specification of the function type signature in the Jluna framework.
  \item \textit{Boxing}: Conversely, Boxing is applied to all return values of the function, which are then passed back to the host language, again as an implicit wrapper around the bridged function.
\end{enumerate}

\begin{figure}[ht!]
  \centering
  \begin{tikzpicture}[scale=1,
                      every node/.style={font=\scriptsize}]

    \begin{class}[text width=2.8cm]{Hittable}{0,3}
      \attribute{+ bbox  : AABB}
      \attribute{+ hit(ray : Ray) : Bool}
    \end{class}

    \begin{class}[text width=2.8cm]{BVHNode}{+2,0}
      \inherit{Hittable}
      \attribute{+ left  : Hittable*}
      \attribute{+ right : Hittable*}
      \attribute{+ bbox  : AABB}
      \attribute{+ hit(ray : Ray) : Bool}
    \end{class}

    \begin{class}[text width=2.8cm]{Shape}{-2,0}
      \inherit{Hittable}
      \attribute{+ material : Material}
      \attribute{+ bbox     : AABB}
      \attribute{+ hit(ray : Ray) : Bool}
    \end{class}
  \end{tikzpicture}
  \caption{Simplified class diagram of the BVH tree structure and its polymorphic components.}
  \label{fig:bvh_class_diagram}
  \Description{A class diagram showing the abstract Hittable type, the BVHNode inheriting from Hittable, and Shape inheriting from Hittable, illustrating the polymorphic structure of the BVH tree.}
\end{figure}

With this in mind, we can already make a few key observations. As stated before, the library is designed to disallow abstract types entirely; hence, this would likely form the first step to our extension. Furthermore, as this library was not designed with abstract types in mind, it naturally did not implement a way to accommodate such types or their attributes in the type implementation (Julia instantiation), boxing, or unboxing. From this, a few interesting questions arise as to how we proceed.

\begin{enumerate}
  \item \textit{How do we encode the abstract base class?} Given that the library was designed to disallow abstract types fully, the type specification needs to implement a mechanism to express subtyping relationships.
  \item \textit{What are the effects of this constraint on the implementation of the library as a whole?} More broadly, this question of extensibility in library design often appears. One interesting aspect of Jluna, specifically, is its heavy usage of compile-time metaprogramming and highly generic type encodings, as opposed to, for example, concrete wrappers or more traditional parsing approaches, as a means of generating bindings. A natural question that then arises is how the extensive application of these paradigms affects the ease of extensibility, given the constraints placed on the initial implementation.
  \item \textit{Unboxing - How do we parse the abstract Julia type in C++?} Naturally, as the nature of our problem is bidirectional, we need to reconcile how the low-level BVH tree representation in Julia is parsed in C++.
  \item \textit{Boxing - How do we represent parse the abstract C++ type to create the corresponding Julia type?} Similarly, we need to ensure that the abstract C++ type is not only implemented in Julia but that we also have a means of moving the C++ type to this corresponding implementation.
  \item \textit{Constraints - How do we loosen constraints while retaining safety?} We need to ensure that we can represent arbitrary BVH trees and establish a method to relax the constraints on type specifications without compromising the system's safety. In other words, we should define a new set of guidelines for what qualifies as a valid polymorphic type specification.
\end{enumerate}

\subsubsection{Metaprogramming in Jluna}

Before elaborating on the actual code that implements this extension, there are two points we find important to establish in more detail before moving on. We have mentioned several times the use of compile-time metaprogramming in the Jluna library. While we have provided a high-level overview of what this entails, we want to dedicate some time to explaining certain aspects of this design paradigm in the original implementation. Additionally, we will examine how these elements were utilized to address some of the questions that emerged during the development of the extension. We begin by providing a brief overview of templates in C++ and how they are used in the Jluna library.

\begin{listing}[ht!]
  \begin{fileminted}{libs/jluna/.src/usertype.inl}{cpp}
template<typename T>
bool Usertype<T>::is_implemented() {
    return _implemented;
}
\end{fileminted}
  \caption{Jluna implemented query function}
  \label{listing:jluna_implemented_func}
\end{listing}

In \ref{listing:jluna_implemented_func}, we see the most basic application of a template, one that should be familiar to most programmers who have used generics before. In this instance, for the set of specified types $T\in \mathbf T$, we instantiate the function for each type in the set. A few observations to make here is that the type $T$ in this isolated instance is not bounded, which is to say that, ignoring \texttt{Usertype}, this function could be instantiated or specialized with any arbitrary single type. To allow for some level of restriction in the types that are permitted to specialize a template, C++ implemented its version of constrained generics called concepts. Concepts define compile-time predicates which express a requirement on a template argument. We can see a basic application of this here.

\begin{listing}[ht!]
  \begin{fileminted}{libs/jluna/include/usertype.hpp}{cpp}
template<is<double> T>
unsafe::Value* box(T value) {
    return jl_box_float64((double) value);
}
\end{fileminted}
  \caption{Jluna concept function}
  \label{listing:jluna_concept_func}
\end{listing}

In this case, we use a concept to call the most appropriate variant of the \texttt{box} function for the double type. Concepts enable static method dispatch; in other words, C++ uses constrained parametric polymorphism to emulate the behavior of dynamic dispatch, as the constraints define a subtyping relationship. Naturally, the language extends this idea to multiple types of arguments, each with its own set of constraints. A question that might arise from seeing this language definition is: What if we want to define a template with a set of types of arbitrary length? Generics that accept an arbitrary set of types are known as variadic generics, also referred to as variadic templates or template parameter packs in the context of C++. Furthermore, the concepts can also be applied to these variadic templates. Once again, we can see an example of this being applied in the Jluna library.

\begin{listing}[ht!]
  \include{code/jluna_variadic_concept_func}
  \caption{Jluna variadic concept function}
  \label{listing:jluna_variadic_concept_func}
\end{listing}

On a high level, all that's happening in \ref{listing:jluna_variadic_concept_func} is that we are checking the predicate that for each member of \texttt{Args\_t}, which represents an arbitrary length set of function arguments, the type of this argument is boxable and hence transferrable across the language boundary. In this snippet, we see two new pieces of notation introduced, namely `\texttt{...}', which is used to denote a variadic generic, and `\texttt{\&\&}', which in this represents an \textit{rvalue reference} which is a reference that binds to an rvalue (temporary value) which is about to be destroyed enabling move semantics for all arguments passed to the function.

A final relevant construct worth mentioning is the template metafunction, a basic example of which can be seen in Listing \ref{listing:template_meta_func}.

We can see a basic example of a particularly relevant construct in template metaprogramming, which enables the above-mentioned functionality. There are various interpretations of this construct, depending on its application; in the most direct case, it's simply a generic struct that exposes its template arguments as statically bound attributes. We can observe an analogous variant of the concept of a generic object in the implementation of traits, where we utilize attributes to define shared behavior for a set of implemented types, much like instantiating a generic type. Traits or interfaces in most languages can then be value-level abstractions, as they provide an abstract or shared view of the behavior exhibited by some set of types. Moving to a higher level of abstraction, we then come to the idea of type traits, which are traits that, as opposed to describing shared behavior, describe a type itself. In the context of C++, this is applied heavily in conjunction with constraints. If we have a means of expressing specific properties of types, it's not too far-fetched to then use those properties to define compile-time constraint predicates. Stated differently, constraints in the context of C++ are type traits that utilize partial specialization to implement compile-time control flow, which in turn enables compile-time predicates on types. The more broad terminology of this idea is a template metafunction, intending to express a function that operations in terms of template arguments, or just types, as opposed to values and `returns' a type or computation applied to this type. For completeness, we should also address the fact that the example \ref{listing:template_meta_func} has a non-type parameter. C++ allows for the passing of constant value-level parameters to templates, which we can use to define compile-time constants that are applied during the template instantiation process. A more formal way of viewing this is that template metafunctions can be defined as dependent types.

\begin{listing}[ht!]
  \begin{fileminted}{src/cpp/example/mfuncex.cpp}{cpp}
template <typename T, const char* name>
struct MetaInfo { 
    using type = T;
    static constexpr const char* name = name;
};
\end{fileminted}
    
  \caption{Template metafunction function}
  \label{listing:template_meta_func}
\end{listing}

\subsubsection{Jluna's Internals}

In the most abstract basic sense, we can think of boxing and unboxing to express the idea of a cross-language assignment. The general goal with binding generators, at least those that focus on very idiomatic language usage, is to provide enough abstractions such that working with two languages feels as if we are working with one. More concretely, the goal is to mitigate, as much as possible within certain constraints, the manual specification of bindings and type marshaling code for handling incompatibilities. Hence, if we go back to our basic idea, what we would then strive for would be something in the following form.

\begin{listing}[ht!]
  \begin{fileminted}{example/assignmentgoal.cpp}{cpp}
MyType my_type_object = foreign_function();
\end{fileminted}
  \caption{Cross language assignment goal}
  \label{listing:assignment_goal}
\end{listing}

If we then compare this to the best analog in Jluna, we have this.

\begin{listing}[ht!]
  \begin{fileminted}{example/jlunaproxyassign.cpp}{cpp}
RGBA cpp_rgba = Main.safe_eval("return RGBA()");
\end{fileminted}
  \caption{Jluna proxy assignment}
  \label{listing:jluna_proxy_assign}
\end{listing}

Here, `\texttt{Main}' represents the module scope, and `\texttt{safe\_eval}' is a wrapper around the Julia C-API function, which uses the runtime to evaluate a snippet of code and return anything specified in a low-level representation. As might already be evident from the code itself, an implicit transformation is occurring. We should clarify that naturally, for this to work, we need existing mirrored-type representation. As we established before, Jluna works quite explicitly by delineating its framework into specification and logic, the specification being the type lens. Furthermore, we know that Julia itself has a variant of lenses, or more simply speaking, it has accessor methods for objects. The main question that arises here what is the implicit transformation happening? The most apparent answer is, well, clear: it's unboxing. If we have an implicit transformation for some low-level representation, then the assignment must be unboxing. While this is accurate, it doesn't address the relevant, more nuanced point of what unboxing and, for that matter, boxing is. 

We can think of it on an intuitive level by using the analogy of types of types as recipes. In this context, we can think of a recipe as having three important parts: the ingredients, the instructions, and the tools. The analogy here is that ingredients represent the abstract properties of our constrained type, while the tools represent how we handle these ingredients. This leaves us with instructions that explain how we use the tools in conjunction with the ingredients to obtain our final product, namely the desired type. Since we have the tools and ingredients, we then know that the instructions must represent the logic of how we use these tools to construct a type. Unboxing and boxing then reduce to the question of, for a language $A$, given the ingredients (our type representation) and the tools (data accessors), how do we use these to construct the type in language $B$ and vice versa. 

Some important takeaways from this are that a type representation, regardless of any explicit specification, will ultimately always be a function of the tools, ingredients, and instructions; i.e., the complexity of the type we can construct in the corresponding language only aligns with the specification insofar as it reflects their fundamental overlap. Secondly, there is a dependency between tools, instructions, and generality of constructions. The more a tool is centered on a specific means of construction, the less it depends on instructions. However, this comes at the cost of losing generality, as the tool is inherently less flexible in its application. Finally, the instructions themselves are a function of the tools and ingredients. In the context of binding generators, this means that the logic of how we construct a type is a function of the data accessors and the type representation. 

\begin{figure}[ht!]
  \begin{tikzpicture}[subgraph text bottom=text centered, subgraph nodes={font=\itshape, black}, nodes=draw]
    \graph [tree layout, grow=right, sibling sep=5mm, level sep=5mm, component align=center] {
        meta constructors -> {type representation, type transformations},
        type transformations -> {type representation, meta constructors},
        meta constructors <-> type transformations
    };
\end{tikzpicture}
  \caption{Basic dependencies of the binding generator}
  \label{fig:bg_basic_dependencies}
  \Description{}
\end{figure}

In \ref{fig:bg_basic_dependencies}, `meta constructors' refer to type construction tools (e.g., lenses), and `type transformations' in this context are simply boxing and unboxing. What follows from this is that if we do want to extend a type of representation, which we do in our case, we must consider a few points.
\begin{enumerate}
  \item \textit{Ingredients}: We have to clearly define what the extended ingredients are for both languages, i.e., what are the new properties of the type representation.
  \item \textit{Tools}: We must ensure that our tools are sufficient for working with the new ingredients, i.e., do we have the necessary data accessors?
  \item \textit{Instructions}: We have to define additional logic for how to use the tools in conjunction with the ingredients to construct the new type.  
\end{enumerate}

With this, we have sufficient knowledge to understand how the original implementation follows from the ideas introduced. We can consider two abstract forms of types: fundamental and grouped. A record represents a grouped type comprising fields that contain either fundamental types or other grouped types. As we cannot trivially bridge a record (as it has unbridged members), we must first bridge the fundamental types before the record. To bridge any type, there needs to exist a context in which we have access to our meta constructors, boxing and unboxing logic, and type representations. Furthermore, we need some way to access this context. We can see the approach taken by the author below. 

\begin{listing}[ht!]
  \begin{fileminted}{jluna/usertype.inl}{cpp}
(...) {
    _mapping.insert({Symbol(name), {
        [get](T& t_inst) -> unsafe::Value* {
            return jluna::box<Field_t>(get(t_inst));
        },
        [set](T& t_inst, unsafe::Value* val) -> void {
            set(t_inst, jluna::unbox<Field_t>(val));
        },
        Type((jl_datatype_t*) jl_eval_string(
            as_julia_type<Field_t>::type_name.c_str())
        )
    }});
}
\end{fileminted}
  \caption{Add property function}
  \label{listing:add_property}
\end{listing}

In Listing \ref{listing:add_property}, we can see how the author relates the type specification to a context in which it can be applied to perform bidirectional mappings of this type. The dictionary `\texttt{\_mapping}' associates the field name `\texttt{symbol}' (wrapped by a Symbol class) with two lambdas which capture the meta constructors (getters and setters) within their context through C++ lambda capture notation `\texttt{[captured](...args) -> ret\_ty \{ \}}', and a wrapped type representation. In this case, we can see that when `\texttt{Field\_t}' is a fundamental type, then the box call will be to a simple boxing method, e.g., for primitive conversions. When `\texttt{Field\_t}' is a record type, i.e., it has itself more fields (assuming it has a UserType instantiation), we call its box method, which in turn will call the box methods on its field. For completeness, we present an adapted pseudocode version of the box and unbox approach of the original implementation in \ref{listing:bridge}.

\begin{listing}[ht!]
  \begin{fileminted}{jluna/usertype.inl}{cpp}
template<typename A>
A Usertype<A>::bridge(B* in) {
    auto out = A();

    for (auto& pair : _mapping)
        std::get<bridge_tool>(pair.tools)(out, 
            jluna::safe_call(
                B_get("accessor"), in,             
                (B*) pair.symbol 
            )
        );
        
    return out;
}
\end{fileminted}
  \caption{Bridge implementation}
  \label{listing:bridge}
\end{listing}

The marshaling code in \ref{listing:bridge} provides a simplified representation of the box and unbox processes. We iterate over the fields specified in the mapping, and for each field, we apply our bridging logic using the meta constructors within the map. These constructors will either bridge the type directly by invoking the fundamental bridging code or, as previously mentioned, call the corresponding bridge method for the record type associated with a specific field, and so on.

A core attribute of the implementation is that it fundamentally relies on nested records existing as fully walkable trees; in other words, the bridging functions must always be supplied with concrete types such that there exists a path over an object that can be fully traced. Hence, this implementation does not account for the case where we encounter an abstract type. The problematic nature, then, is that our implementation of the BVH tree is fundamentally based on abstraction for its representation, or more concretely, as stated before, we have polymorphic nodes represented by an abstract class. The aim of the extension is, first and foremost, to determine how we can instruct this process on what to do if it encounters an abstract type. 

\subsection{Implemented Extension}

Let's start by creating a high-level breakdown of the user-facing and internal flow of the library, as shown in \ref{fig:jluna_stages}. We could linearly tackle this problem, although it may become somewhat problematic since we might define a specification that is misaligned with the actual type mirroring process. The issue here stems from the different dependencies in the binding generator, as shown in \ref{fig:bg_basic_dependencies}. 

\begin{figure}[ht!]
  \centering
  \begin{tikzpicture}[subgraph text bottom=text centered, subgraph nodes={font=\itshape, black}, nodes=draw]
  \graph [layered layout, grow=down, sibling sep=5mm, level sep=5mm, component align=center, component direction=270] {
    // [component align=center] {
      "type specification" -> [dashed] {
        "type implementation" -> "type generation", 
      }
    };

    "type mirroring"[draw] // [layered layout, grow=down] {
      "type implementation",
      "type generation"
    },

    "type specification" -> "type registration",

    "type mirroring" ->[dashed] "type registration",
  };
\end{tikzpicture}
  \caption{Stages of Jluna's type binding generation process; $\dashrightarrow$ represents user-facing flow, i.e., how the user interacts with the library; $\rightarrow$ represents the internal flow of the library}
  \label{fig:jluna_stages}
  \Description{}
\end{figure}

\subsubsection{Representing Fields \& Types}

We can approach things in a top-down fashion, starting from the fundamental type representation and progressing to the actual bridging code. To bring us back to where we started with our earlier developed formalisms, we had the most basic form of our record represented as a type family indexed by a set of labels. We had three main components to this: the labels, the object type, and the dependent field type. The corresponding implementation in C++ can be seen in Listing \ref{listing:type_meta_func}.

\begin{listing}[ht!]
  \begin{fileminted}{src/cpp/usertype.cpp}{cpp}
template <
    ct_string Name, 
    typename T, typename M, 
    M T::* Member
>
struct Lens<Name, Member> {
    using type  = internal_t<M>;
    
    static M T::* member_ = Member;

    static auto label() { 
        return Name; 
    }
    static auto get(T& obj) -> M { 
        return obj.*member_; 
    }
    template <typename U>
    static auto set(T& obj, U&& value) -> void {
        obj.*member_ = value;
    }
};
\end{fileminted}
  \caption{(Template) Field Type meta-function}
  \label{listing:type_meta_func}
\end{listing}

As a reminder, unlike a typical function, a template metafunction takes a set of statically bound types as its arguments and returns a type as its result. In our case, we have two main components: the name of the field and the \textit{pointer-to-member} type. The name is associated with the Lens using a \texttt{ct\_string} type, which is a struct that holds a compile-time string and provides added modularity when declaring lenses, as opposed to explicitly declaring static strings. The pointer-to-member type is a commonly used C++ construct that enables us to reference a member of a class or struct. A pointer-to-member type is defined as `\texttt{M T::* name}', where `\texttt{M}' is the type of the member, `\texttt{T}' is the type of the class or struct, and `\texttt{name}' is the name of the member. Internally, we expose a stripped \texttt{type} type, which represents the parameterized type variable in case the member is defined as being specialized over a type. To fulfill the requirements of our Lens, we also implement the generic getter and setter methods, in addition to a label reflection mechanism to access the field name from a static context. A notable point of this Lens is that with the setter method, we introduce an additional type variable $U$ to account for the case of assigning a derived pointer to a base-type member. 

To contrast this with the original implementation here, instead of binding the field properties to the context of the \texttt{add\_property} function, we isolate them as an entirely static metafunction. There are several benefits to this approach, but the primary motivation is to create a clear separation and generalization between how we specify a type, the tools used to construct it, and the actual logic that applies these tools. Another benefit of this is that, while currently, there exists no native reflection within C++, it should be a relatively non-trivial task to generate these metafunctions at compile time once reflection is introduced in C++26.

A whole type, or more concretely, in our case, object, is then a composition of these fields. As metafunctions typically act on the type level rather than the value level, we pass them to functions as template arguments. Here, we have variadic templates come into play, which allows us to pass an arbitrary number of our metafunctions to a function, in turn allowing us to construct this composition of fields. The original implementation divided the process of defining types into two parts: the C++-sided definition, i.e., type registration through the \texttt{add\_property} function, and the Julia-sided definition, i.e., type implementation and mirroring. As we wanted to generalize the field or property definition as a function that registers a type over a group of fields as opposed to for each field individually, we adapted the logic into the \texttt{init\_type} function, the signature of which can be seen in listing \ref{listing:init_type_signature}.

\begin{listing}[ht!]
  \begin{fileminted}{libs/jluna/include/usertype.hpp}{cpp}
template <typename... FLens, typename... DT>
static void init_type(
    TL<FLens...>, 
    TL<DT...> = TL<>()
);

template <typename... FLens, typename... DT>
static void init_dmap(TL<FLens...>, TL<DT...>);
\end{fileminted}
  \caption{Type initialization signature}
  \label{listing:init_type_signature}
\end{listing}

One caveat we encounter when working with variadic functions is how to delineate between distinct sets of arguments. In our case, we needed not only to pass the set of fields as $P$ but also to pass the derived types relevant for a particular field to perform the necessary concrete cast. To separate the two sets of arguments, we define a variadic container as being a metafunction that acts as a means of enclosing a set of types. Here, we pass these lists as value-level metafunctions, which then allows us to define two distinct sets of type lists \texttt{FLens} and \texttt{DT}, with \texttt{DT} representing our derived types. At compile time, this serves as a means to specialize our type registration for the given specification. At runtime, this specialization is then used to register all the types for the C++-sided code. 

The Julia-sided type mirroring happens in C++ through the implementation declaration. We can see the signature of this in the listing \ref{listing:type_impl_signature}. The most relevant parts here are the two type parameters $T$ referring to the type of the object and $U$ being the supertype. On a high level, we can think of this as expressing to Julia that within the context of a module (the default just being the global scope), we want to establish the relationship $T <: U$. The type generation triggered by this call is then executed at runtime and calls into the Julia-side portion of the library actually to create a type with this relationship. To avoid redefinition errors, check if a type has already been defined on the Julia side within the specified module scope and reuse existing definitions if they match our C++-sided specification. 

\begin{listing}[ht!]
  \begin{fileminted}{libs/jluna/include/usertype.hpp}{cpp}
template<typename T>
template<typename U>
void Usertype<T>::implement(unsafe::Module* module)
\end{fileminted}
  \caption{Type implementation signature}
  \label{listing:type_impl_signature}
\end{listing}

How we then use both this C++-sided initialization and its corresponding runtime registration, as well as the Julia-sided mirroring, is shown in an example with listing \ref{listing:usertype_ex_appl}. Here, \texttt{TL} represents our basic type list, which contains the properties of the sphere type; the second argument is the relevant derived types for these properties, in this instance, the Lambertian and Metal materials. The implementation then expresses the relationship $\text{Sphere} <: \text{Hittable}$ and, at runtime, uses the information from the registration to mirror the type with the given properties.

\begin{listing}[ht!]
  \begin{fileminted}{src/cpp/usertype.cpp}{cpp}
Usertype<Sphere>::initialize_type(
    lenses<
        Lens<"center"_ct,  &Sphere::center>,
        Lens<"radius"_ct,  &Sphere::radius>,
        Lens<"r_squared"_ct,&Sphere::r_squared>,
        Lens<"mat"_ct,     &Sphere::mat>,
        Lens<"bbox"_ct,    &Sphere::bbox>
    >(),
    types<Lambertian, Metal>()
);

Usertype<Sphere>::implement<Hittable>();
\end{fileminted}
  \caption{User type example application}
  \label{listing:usertype_ex_appl}
\end{listing}

\subsubsection{Semantics of Unboxing \& Boxing}

Starting with the unboxing process, we can remind ourselves that the core issue here was that the unboxing mechanism we wanted to use for a field of an abstract type had to correspond to the underlying concrete type, given the library's general design and resulting constraints. The key aspect of the solution in the unboxing case is that, since Julia uses parametric structural subtyping, we typically know the concrete type from the Julia side by simply querying the type info attribute of the low-level object representation used within the C++-sided Jluna code. From this principle, we can construct the concrete unboxing mechanism within Jluna by establishing a correspondence between a pair consisting of a field label and a valid subtype of this field and the unboxing mechanism, i.e., the setter function for this field and its subtype.  

\begin{listing}[ht!]
  \begin{fileminted}{libs/jluna/include/usertype.hpp}{cpp}
static inline std::unordered_map<
    std::pair<std::string, std::string>, 
    std::function<void(T&, unsafe::Value*)>
> _type_dispatch_map = {}; 
\end{fileminted}
            
  \caption{Type dispatch map}
  \label{listing:type_dispatch_map}
\end{listing}

The resulting unboxing mechanism is essentially a simple lookup into this map. The actual instantiation of this map then has two main cases: fields with concrete types and fields with abstract types. If a field has a concrete type, we create a trivial mapping from the concrete type to the concrete setter. If the field has an abstract type, we resolve all valid derived types and then create a set of mappings to all valid derived types. In simple terms, we create a map consisting of all valid concrete instantiations for all fields of an object. We can define this by cases as follows. 

\begin{definition}[Dispatch Map Entries]
  \[
\text{entry}( l:T_{l}) =\begin{cases}
\langle l,\ T_{l} \rangle  & \texttt{if isConcrete}( T_{l})\\
\prod _{i:I} \langle l,R_{i} \rangle \  & \text{for all }( R_{i})_{i:I} < :T_{l}
\end{cases}
\]

  The entry of a dispatch map for a record of type $T$, a label $l$ and a member type $T_l$ is defined as the projection $\langle l, T_l\rangle \rightarrow \text{set}\langle T_l\rangle(x)$ when $T_l$ is concrete, and the projection $\prod_{i:I}\langle l, R_i \rangle \rightarrow \text{set}\langle R_i\rangle(x)$ when $T_l$ is abstract and $R_i$ are the valid derived types of $T_l$. 
\end{definition}

One notable feature that plays a role in dispatch map instantiation is the type filtering mechanism, the pseudo-code of which is shown in Listing \ref{listing:type_filter}. An early design choice was that in the type specification, all derived types relevant to a particular object would be passed as the aforementioned template pack, a benefit of this being that with manual specification, we avoid forcing repeated specifications, so if two fields, for example, have the same abstract type we only need to specify the associated derived types once. Furthermore, we reduce implementation complexity somewhat as the set of types does not exist as a member of a Lens but instead as just an auxiliary pack. A key downside to this approach is that we naturally lack a direct reference to the fields that Lens has and the relevant derived types. Thus, when iterating over the fields, we must first filter out derived types that are irrelevant to a specific field to avoid instantiating invalid setters. In simple terms, if a field, for example, has a concrete type $A$, then we don't want to instantiate setters for any derived types irrelevant to $A$. 

\begin{listing}[ht!]
  \begin{fileminted}{libs/jluna/.src/usertype.inl}{cpp}
for_each<FilterType<
    FieldType, is_base_of, DerivedType...
>>();
\end{fileminted}
                
  \caption{Type filter}
  \label{listing:type_filter}
\end{listing}

The mechanism is generally straightforward. It consists of an initial type $T$, a binary type predicate $P : \langle A, B \rangle \rightarrow \texttt{Bool}$ and a set of types $\mathbf T$, the resulting filtered type set $\mathbf T_f$ is then the set of types $\mathbf T$ for which the predicate $P(T_1, T_2)$ is true. There are various points to be made regarding the pros and cons of this approach. The primary motivation in this case was to simplify the type specification and delegate the actual filtering process to the compiler. Let's consider the automated case of parsing types. This mechanism might have less utility as we could implicitly construct the valid type specification as a part of the parsing algorithm itself. However, it still increases the complexity of the implementation somewhat, regardless, as an implicit valid specification would mean the implementation does not act linearly but instead on a tree of types from the outset. 

\begin{listing}[ht!]
  \begin{mathalgo}{Unboxing usertype}
foreach_indexed<FLens...>([&]<Lens FL, size_t I> { 
    invoke_with<FType<FL>>(
        match {
            [&]<is_jlbits FT>(Value* v) { 
                FL::set(out, as_pod<FT>(v)); 
            }, 
            [&]<typename FT>(Value* v) {
                dispatch_tptr<FT>(out, v);
            }
        }, 
        nth_field<FL, I>(in /* Value* in */)  
    );          
});
\end{mathalgo}
  \caption{Pseudocode snippet of the unboxing process}
  \label{listing:unbox_usertype}
\end{listing}

Two important optimizations we use to accelerate the unboxing process are immediately setting \textit{piece of data} (PoD) fields and using a \textit{polymorphic inline cache} (PIC) in conjunction with the dispatch map. The first optimization leverages the fact that if a type we receive from Julia has a memory layout that is compatible with the corresponding C++ type, then we can reinterpret the memory payload of the Julia object as the C++ type and set the field directly. This is particularly useful in our case for static arrays as the memory layout of a Julia-sided `\texttt{SVector{Float64, 3}}' is identical to a C++-sided `\texttt{std::array<double, 3>}'. The second optimization is the use of a PIC, a common optimization technique in dynamic languages that speeds up method dispatch. In our case, we use the PIC to cache the results of the setter function lookups in the dispatch map. This means that if we have already looked up the setter function for a particular field and concrete type, we can reuse the cached result, which is especially useful for frequent, similar lookups. The size we chose for our PIC is four entries before we resort back to a regular lookup in the dispatch map and evict entries. Further testing is required to evaluate the broader implications of this choice. However, it falls within the original empirically determined \cite{HlzlePICs} and commonly implemented range of around 4-10 entries (e.g., 4 for V8 \cite{v8_wasm_constants}, 6 for SpiderMonkey \cite{mozilla_icstate}). 

Similarly, in the case of unboxing, the boxing process requires that we know the underlying concrete type so we can use the correct mechanism to map to the Julia-side type. We utilize the exact type pattern matching applied in the dispatch map construction to differentiate between boxing concrete and abstract types. The concrete case is a direct call to the appropriate boxing function, essentially unchanged from the original implementation. The abstract case requires us to determine the correct concrete type. Unlike in the case of Julia, C++ provides no clean, native way of querying the concrete type of an abstract pointer. The most trivial and safe solution, albeit at the cost of performance, is to brute-force test all valid derived types using a dynamic cast, essentially treating the abstract type as a union. As the library is designed to require no additional modifications to types themselves, we cannot expect any self-tagging mechanisms to be present. Hence, any scheme for determining the concrete type must use only the information provided by the type specification and the resulting runtime types we are passed. 

\begin{listing}[ht!]
  \include{code/boxing_simple}
  \caption{Pseudocode snippet of the boxing process}
  \label{listing:box_usertype}
\end{listing}

The fast, albeit dangerous, solution we ended up implementing to avoid brute force dynamic casts was to create a dispatch map based on the vptr of the derived types. The fact we can exploit's that, despite there being no strict specification of how vptrs are implemented, the convention most compilers end up adopting is having the vptr as a hidden 0th member of the class. Furthermore, since vptrs definitionally act as unique tags to a type-specific table of virtual functions, we can leverage this to create a mapping of vptrs to unbox functions of the corresponding type. In essence, we are creating an auxiliary extension to the vtables of all derived types using the vptrs as a key to the correct boxing function. One future optimization is to make our PIC bidirectional by generalizing it to support both boxing and unboxing.

\subsection{Evaluation}

Our evaluation will consist of two main parts: the first is a performance evaluation of the library, and the second is a qualitative evaluation of the developer and user experience. The performance evaluation will detail the benchmarking setup, the performance of boxing and unboxing, and the overhead introduced by the application of JuliaCall and Jluna.  

\subsubsection{Benchmarking \& Performance}

On a high level, both JuliaCall and Jluna function on similar principles for establishing a line of communication between C++ and Julia. Both libraries provide us with a means to start the Julia runtime, call Julia-sided functions, or conversely define C++ or Python-sided functions that can be called from Julia. In our applied instance, we want to intercept the call to the BVH traversal function and then call into our foreign functions from Julia. The approach we take to intercept the Julia-sided call is to redefine the traversal function within the scope of the Julia ray-tracing module, replacing the original function with one that calls into our foreign code. The internal trampoline used by the redefined function is created using the respective interop libraries, both of which, in essence, more or less generate wrappers around Julia's `\texttt{ccall}' function. 

\begin{figure}[ht!]
  \centering
  \begin{sequencediagram}
    \newinst {img}{:Image Gen}
    \newinst  {proxy}{:Proxy}
    \newthread  {api}{:Interop Lib}
    \newinst  {bvh_p}{:BVH}

    \begin{call}{api}{render(scene)}{img}{}
        \begin{sdblock}{proxy module}{} 
        \begin{sdblock}{draw}{pixel $\times$ sample}
            \begin{call}{img}{hit(args…)}{proxy}{HitInfo}
                \begin{call}{proxy}{hit(args…)}{api}{HitInfo}
                    \begin{callself}{api}{unbox(args)…}{…}{}
                    \end{callself}
                    
                    \begin{call}{api}{hit(…)}{bvh_p}{HitInfo}
                        \begin{callself}{bvh_p}{hit(…)}{…}{}
                        \end{callself}
                    \end{call}

                    \begin{callself}{api}{box(HitInfo)}{…}{}
                    \end{callself}
                \end{call}
                    
                \node[anchor=east, xshift=-4pt] (t0) at (cf3) {$t_0$};
                \node[anchor=west, xshift=4pt] (t0) at (ct4) {$t_1$};
                
                \node[anchor=west, xshift=4pt] (t1) at (rt4) {$t_2$};
                \node[anchor=east, xshift=-4pt] (t1) at (rf3) {$t_3$};       
                \draw [decorate, 
                    decoration = {brace,
                    raise=5pt,
                    mirror,
                    amplitude=3pt}]%
                    let
                        \p1=(ct3),
                        \p2=(cf4)
                    in
                        (\x2,\y1) -- (\p2)
                    node[pos=0.5,left=7pt]{$\Delta t_{u}$};

                \draw [decorate,
                    decoration = {brace,
                    raise=5pt,
                    mirror,
                    amplitude=3pt}]%
                    let
                        \p1=(cf4),  
                        \p2=(rf4)
                    in
                        (\x2,\y1) -- (\p2)
                    node[pos=0.5,left=7pt]{$\Delta t_{r}$};   
                    
                \draw [decorate,
                    decoration = {brace,
                    raise=5pt,
                    mirror,
                    amplitude=3pt}]%
                    let
                        \p1=(rf4),  
                        \p2=(rt3)
                    in
                        (\x2,\y1) -- (\p2)
                    node[pos=0.5,left=7pt]{$\Delta t_{b}$};
                    
                    \postlevel
                \end{call}
            \end{sdblock}
        \end{sdblock}
    \end{call}
\end{sequencediagram}
  \caption{Benchmarking setup}
  \label{fig:bechmarking_setup}
  \Description{}
\end{figure}

We have created a sequence diagram for the benchmarking setup in \ref{fig:bechmarking_setup}. The first timing point is recorded right before calling the generated ccall trampoline and is denoted by $t_0$. This calls into the respective interop libraries, which proxy the user-defined code in our foreign language, i.e., our rewritten BVH traversals. The second timing point, $t_1$, is recorded once we have entered the user-defined traversal functions, marking the start of the actual work. Similarly, we generate a third timing point, $t_2$, prior to boxing once our traversal is completed and a fourth timing point, $t_3$, after the boxing has concluded and our value is returned to the Julia runtime. Additionally, we also recorded the time spent within Julia's GC for the entire duration between $t_0$ and $t_3$. 

All benchmark data are derived from the task of rendering the same cottage scene using the C++ and Python implementations of the BVH traversal. Naturally, as our focus is on the performance of the libraries rather than the performance of the actual ray-tracing algorithms, we primarily discuss $\Delta t_u$ and $\Delta t_b$, the time spent on unboxing and boxing, respectively. The actual work time, $\Delta t_r$ does have relevance, though, insofar as it relates to the object representations derived from the unboxing process. 

A simple overview of the results can be seen in figure \ref{fig:boxing_unboxing_basic}. The most apparent observation we can make from this is that boxing and unboxing times are considerably higher in C++ than in Python. Taking unboxing as an example, within Jluna, we average around 75.8 µs per unboxing, regardless of all arguments, whereas in Python, we only average around 5.8 µs per unboxing. The boxing times are similarly skewed, with C++ averaging around 59.3 µs per boxing and Python averaging around 3.1 µs per boxing. We can attribute this difference in large part to the nature of the underlying implementations, namely copying vs. non-copying interoperability. 

JuliaCall, like other Python interop libraries, only builds wrapper objects, choosing not to copy the actual data. Instead, it relies on proxying or intercepting getter and setter operations using the Julia C-API. In practice, this means that the only data we create when "unboxing" in JuliaCall are minimal proxy objects, which more or less wrap a Julia-sided pointer. In contrast, Jluna (in the way we apply it) copies over the actual data, as we intended to leave our original C++ implementation unchanged. This means that our unboxing process in Jluna fundamentally necessitated a de facto deep copy of the data, in turn leading to higher unboxing and boxing times. 

\begin{figure}[ht!]
  \centering
  \includesvg[width=0.45\textwidth]{figures/boxing_unboxing_comparison}
  \caption{Boxing and unboxing comparisons}
  \label{fig:boxing_unboxing_basic}
  \Description{}
\end{figure}

We can observe a more granular breakdown of the boxing and unboxing times in figure \ref{fig:boxing_unboxing_mirrored}. An interesting observation from this figure is especially the structural differences in the execution time of different marshaling operations for the ray casts. Within Python, we observe consistent spikes in both boxing and unboxing times, particularly between the 400th and 800th ray casts. It's difficult to say precisely what the source of these spikes is, but a potential explanation could be the Python GC. The C++ implementation behaves considerably more inconsistently; the most striking observation is the sudden high-frequency spikes in the boxing times. We can notice similar, albeit less pronounced, spikes in the unboxing times. When we consider the otherwise relatively flat profile of Boxing and unboxing times in C++, we could attribute inconsistent cache locality as a potential source of these spikes.

\begin{figure}[ht!]
  \centering
  \includesvg[width=0.45\textwidth]{figures/boxing_unboxing_mirrored.svg}
  \caption{Boxing and unboxing mirrored}
  \label{fig:boxing_unboxing_mirrored}
  \Description{}
\end{figure}

One observation made while executing the Python implementation was that there were occasional stutters in the execution. However, this appeared to contradict the otherwise relatively consistent performance noted in both \ref{fig:boxing_unboxing_basic} and \ref{fig:boxing_unboxing_mirrored}. To investigate this further, we looked at specifically the time spent in the Julia GC within the Python implementation. It's clear, looking at \ref{fig:boxing_unboxing_mirrored}, that these stutters don't appear to occur in the boxing or unboxing times; hence, they must be occurring during the execution of the actual work within Python. To investigate this more clearly, we graphed the total overhead of the Python marshaling process, along with the time spent in the Julia GC for each ray cast in figure \ref{fig:overhead_mirrored}.

\begin{figure}[ht!]
  \centering
  \includesvg[width=0.45\textwidth]{figures/overhead_mirrored.svg}
  \caption{Overhead mirrored}
  \label{fig:overhead_mirrored}
  \Description{}
\end{figure}

Something that stands out in \ref{fig:overhead_mirrored} is just how significantly larger the GC spikes are during the execution of the Python BVH traversal than the actual overhead of the marshaling process. In other words, the overhead of the marshaling process in Python is minimal, with no recorded time spent within the Julia GC. However, the actual application of the wrapper types within our Python BVH traversal contributes significantly to the overall runtime. We can emphasize this more clearly by looking at the cumulative overhead in figure \ref{fig:cumulative_overhead}.  

\begin{figure}[ht!]
  \centering
  \includesvg[width=0.45\textwidth]{figures/cumulative_overhead.svg}
  \caption{Cumulative overhead}
  \label{fig:cumulative_overhead}
  \Description{}
\end{figure}

In \ref{fig:cumulative_overhead} we can see that the marshaling overhead shows us the same profile we observed in \ref{fig:boxing_unboxing_basic}, but the overhead introduced by the time spent in the GC for the BVH traversal of the Python implementation puts us at times almost on the same level as the overhead introduced by just copying the actual data as opposed to holding the pointers and proxying all operations. 

\begin{table}[ht!]
  \centering
  \begin{tabular}{lrrrr}
    \toprule
 Metric   & \multicolumn{2}{c}{Python}       & \multicolumn{2}{c}{C++}          \\
             & Mean (µs) & $\sum$ (µs)      & Mean (µs) & $\sum$ (µs)      \\
    \midrule
 Unboxing & 5.8284       & 8\,346.32       & 75.7565      & 105\,983.308    \\
 Boxing   & 3.1438       & 4\,501.9580     & 59.3206      & 82\,989.486     \\
 Actual   & 1\,874.7250  & 2\,684\,606.196 & 0.1723       & 241.0350        \\
 Overhead & 8.9723       & 12\,848.278     & 135.0771     & 188\,972.794    \\
    \bottomrule
  \end{tabular}
  \caption{Comparison of average and total times for unboxing, boxing, actual work, and overhead in Python vs.\ C++.}
  \label{tab:perf-metrics}
\end{table}

Something we can take away from the performance evaluation is that while the C++ implementation does have a higher overhead in terms of boxing and unboxing, as seen again in \ref{tab:perf-metrics}, it does not suffer from the same GC overhead as the Python implementation. This is primarily because we copy over the actual data in C++. As the Julia GC only executes within the context of the Julia runtime, it's clear that in this instance, a drawback of repeatedly calling into the Julia runtime, as is the case with JuliaCall, is that we incur a significant overhead in terms of the time spent in the GC. In contrast, Jluna has no recorded time spent in the Julia GC, as the user-defined code is executed entirely within C++, and the marshaling process is implemented to minimize GC involvement. 

\subsubsection{Developer \& User Experience}

One aspect the original implementation had in mind, which we tried to continue with our extension, was utilizing modern C++ features to improve the developer experience of working with the code. The main aspects of this were the use of well-documented and type-safe template metaprogramming, in addition to functional programming and, more broadly, declarative programming paradigms. In practice, this meant we aimed to design the implementation—especially the more abstract portions of the code—so that they would reflect the most critical parts of the implementation with as little noise as possible. As templates are somewhat notoriously complex to read, we followed the style of the original implementation by keeping heavy template or type predicate code separate from the main implementation, and we ensured that types have clear and descriptive names. Functional or more complex compile-time type constructs were integrated in a way that reflects a sort of two-level language, where abstractions of templates are used in a manner more naturally akin to a functional compile-time programming language, separate from the runtime implementation, and with the underlying type logic being separated from the actual implementation.

The user experience is multifaceted, depending on the user's priorities. While the library does not yet achieve fully implicit interoperability, it does provide a relatively simple interface for declaring types and functions. In its current state, this is a mundane and time-consuming action, depending on the size of the codebase, but it could be automated with relative ease once reflection is introduced in C++26. A key benefit of this library is that, despite requiring some explicit declarations, these are largely separate from the main implementation. This isolation allows for more versatility in the application of the library, as opposed to those that require a tighter coupling between the interoperability code and the main implementations in either language. As mentioned earlier, performance overhead is a factor that requires careful consideration when using the library. To aid in this, a basic user-facing benchmarking tool was implemented for the C++ side, and Julia, being a dynamic high-level language, has various tools to benchmark from its end.

\subsection{Limitations \& Threats to Validity}

We can understand the limitations of this extension and, more broadly, the idea of expanding interoperability libraries to support more complex types by first assessing some important language-level considerations and then examining the broader landscape of the evolution of programming languages and their real-world applications.

\subsubsection{Slow and varying generics support}

Generics support, or more complex type-level programming, is a key aspect of building more implicit, idiomatic interoperability libraries that aren't reliant on external scripts or tools to parse and generate the necessary code. C++ and Julia are both languages that have, in their ways, developed considerably complex type systems. A benefit of this focus in their design is that we can leverage these systems, especially on the C++-sided end, to approach the problem of handling many different types in highly generic ways, thereby minimizing the amount of boilerplate code we need to write. Treating the problem of interoperability with this type-level approach, while quite interesting in this case, could pose a significantly larger challenge in languages that don't have as developed type systems. Furthermore, it's not even a guarantee that languages will continue to develop increasingly complex type systems, as even general-purpose programming languages often end up having domain-specific requirements or design choices that can deliberately limit a highly expressive type system. A type-level approach is a suitable approach in languages that easily support it, but it could be detrimental in languages based on incompatible design philosophies.

\subsubsection{Asymmetric metaprogramming features}

As briefly touched upon in the previous section, due to the varying development and design of different programming languages, they often have very different metaprogramming features. As our implementation is heavily reliant on metaprogramming features for things like type-level functional logic, reflection, type generation, specialization, and so on, similarly to the previous point, it may not be easy to use these design choices as a more general framework for interoperability systems. Metaprogramming exists to various degrees in many languages. However, its feature set should be closely evaluated to determine if it's sufficient to use as the basis for an interoperability library or if other approaches would be more suitable.

\subsubsection{Decentralized language evolution}

If we examine the broader landscape of programming languages, we can observe a general trend: languages tend to evolve within their ecosystems. Naturally, there are exceptions to this trend, and various attempts have been made to build languages and more abstract frameworks that enable different languages to coexist within the same ecosystem. A key example of this is the JVM, CLR, and, more recently, the GraalVM. However, despite these developments, many mainstream popular languages continue to evolve within their respective communities and ecosystems. This is not to say that these languages are entirely isolated from one another. We do often see theoretical or practical aspects of their implementation being inspired by other languages or built on existing tools, a prime example of this being the set of languages developed on LLVM. Even with the often theoretical or technological overlap of languages, a limiting factor in creating traditional interop libraries that enable the bridging of more complex properties of languages, especially ones that diverge in their design philosophies, is that the languages are often inherently not designed to interoperate as we often see them created and evolve to address either a deficit in another language or some completely different domain-specific niche. We have seen various developments aimed at bridging separate language ecosystems, such as WASM or the C-ABI. However, these are often limited, and as we established, they come with certain fundamental drawbacks. We have also observed that languages built on an inherently unifying framework can overcome many of the deficits of more traditional interop libraries. However, existing, well-established language ecosystems continue to be the primary means by which people interact and work with these languages.  

\subsubsection{Theoretical exploration vs practical needs}

A final point to consider is the fundamental reality or question of "\textit{is interoperability past basic constructs even necessary in the first place?}". In an industry setting, especially at larger scales, it's somewhat natural for the environment to focus on the more short-term to mid-term outcomes of a piece of software. This tends to lead to a preference for ease of use, safety, reliability, and other factors. An observation that can be made from this preference for safe, reliable, and easy software in a practical setting is that it often rarely attempts to evolve beyond what is actually needed and what can be reliably trusted. Legacy systems and codebases often result from this preference. While this is not inherently a bad thing, it does, in many instances, mean that, rather than spending much time exploring or ruminating on the latest features in a language, the focus is more on simply working with what we have to get the job done. Exploration and being more open to new technologies tend to be more common in a research setting or sometimes only possible in instances where something is financially viable. Creating more expressive interoperability, while interesting to explore in a theoretical sense, is not necessarily a real practical need for many developers. Research and academia being ahead of what is used in day-to-day life is not a new phenomenon; however, in this instance, there is a clear need for more research to assess whether this approach to interoperability is more broadly applicable, given the constraints elaborated on above.
